% Figure: two paths between the same two states on a p-V plane, illustrating
% that DeltaU is path-independent but W (area) differs.
\begin{figure}[htbp]
\centering
\begin{tikzpicture}
\begin{axis}[
  width=11cm, height=7.5cm,
  xlabel={volume $V$}, ylabel={pressure $p$},
  xmin=0.5, xmax=4.5, ymin=0, ymax=5,
  axis lines=left,
  legend pos=north east,
  legend cell align=left,
  tick label style={font=\scriptsize},
  label style={font=\small},
  legend style={font=\scriptsize},
  clip=false,
]
% State A = (1,4), State B = (4,1)
% Path I: isothermal pV=4 from A to B
\addplot[engblue, very thick, smooth, domain=1:4, samples=50] {4/x};
\addlegendentry{path I (isothermal)}
% Path II: down then across (isochoric then isobaric)
\addplot[engred, very thick] coordinates {(1,4) (1,1) (4,1)};
\addlegendentry{path II (isochoric + isobaric)}
% state markers
\filldraw[engblack] (axis cs:1,4) circle [radius=2.2pt];
\node[font=\scriptsize, anchor=south west] at (axis cs:1,4) {state $A$};
\filldraw[engblack] (axis cs:4,1) circle [radius=2.2pt];
\node[font=\scriptsize, anchor=south west] at (axis cs:4,1) {state $B$};
\end{axis}
\end{tikzpicture}
\caption[State functions versus path functions]{Two processes carry the system
from state $A$ to state $B$. The change in any state function, $\Delta U$,
$\Delta H$, $\Delta S$, is identical for both paths because it depends only on
the endpoints. The work $W = -\int p\,dV$ is the area under each curve, and the
two areas differ, so $W$ depends on the path taken. Heat $Q = \Delta U - W$
then also differs between the paths. Work and heat are path functions; the
state functions are not. This is why a cyclic engine can deliver net work even
though every state function returns to its starting value each cycle.}
\label{fig:vol04ch01:state-vs-path}
\end{figure}
